% 阴阳四象图
% 用途：展示阴阳到四象的演化：阴爻、阳爻 → 少阳、老阳、少阴、老阴，搭配一天和四季
\documentclass[tikz,border=10pt]{standalone}
\usepackage{ctex}
\usepackage{xcolor}
\pagecolor{white}
\usetikzlibrary{arrows.meta, positioning, calc, decorations.pathmorphing, backgrounds}

\begin{document}
\begin{tikzpicture}[scale=1.05]

  % ====== 两仪 ======
  \node[font=\LARGE\bfseries] at (0,7.0) {两仪};

  % 阳爻
  \draw[very thick] (-1.8,5.8) -- (1.8,5.8);
  \node[font=\large\bfseries] at (0,5.3) {阳爻};
  \node[font=\small] at (0,4.9) {显现、主动、明亮};

  % 阴爻
  \draw[very thick] (-1.8,4.1) -- (-0.3,4.1);
  \draw[very thick] (0.3,4.1) -- (1.8,4.1);
  \node[font=\large\bfseries] at (0,3.6) {阴爻};
  \node[font=\small] at (0,3.2) {隐藏、安静、幽暗};

  % 分割线
  \draw[dashed, gray!50] (-2.8,2.7) -- (2.8,2.7);

  % ====== 四象 ======
  \node[font=\LARGE\bfseries] at (0,2.3) {四象};

  % 少阳
  \node[font=\large\bfseries, orange!70!black] at (-5.2,1.0) {少阳};
  \draw[very thick] (-6.5,0.3) -- (-5.0,0.3);
  \draw[very thick] (-6.5,-0.2) -- (-5.5,-0.2);
  \draw[very thick] (-5.1,-0.2) -- (-5.0,-0.2);
  \node[font=\small] at (-5.75,1.7) {阳刚开始生};
  \node[font=\small] at (-5.75,-0.7) {清晨 · 春};

  % 老阳
  \node[font=\large\bfseries, red!70!black] at (-1.7,1.0) {老阳};
  \draw[very thick] (-3.0,0.3) -- (-1.6,0.3);
  \draw[very thick] (-3.0,-0.2) -- (-1.6,-0.2);
  \node[font=\small] at (-2.3,1.7) {阳达到极盛};
  \node[font=\small] at (-2.3,-0.7) {正午 · 夏};

  % 少阴
  \node[font=\large\bfseries, blue!60!black] at (1.7,1.0) {少阴};
  \draw[very thick] (0.4,0.3) -- (1.8,0.3);
  \draw[very thick] (1.1,0.3) -- (1.1,-0.2);
  \draw[very thick] (0.4,-0.2) -- (0.8,-0.2);
  \draw[very thick] (1.4,-0.2) -- (1.8,-0.2);
  \node[font=\small] at (1.1,1.7) {阴刚开始生};
  \node[font=\small] at (1.1,-0.7) {傍晚 · 秋};

  % 老阴
  \node[font=\large\bfseries, purple!70!black] at (5.2,1.0) {老阴};
  \draw[very thick] (4.0,-0.2) -- (4.8,-0.2);
  \draw[very thick] (5.2,-0.2) -- (6.0,-0.2);
  \draw[very thick] (4.0,0.3) -- (4.8,0.3);
  \draw[very thick] (5.2,0.3) -- (6.0,0.3);
  \node[font=\small] at (5.0,1.7) {阴达到极盛};
  \node[font=\small] at (5.0,-0.7) {午夜 · 冬};

  % 底部循环
  \node[font=\normalsize, gray!70!black] at (0,-1.5) {少阳 → 老阳 → 少阴 → 老阴 → 少阳 …};
  \node[font=\small, gray] at (0,-2.0) {阴阳消长，循环不息};

  % 白色背景，防止 GitHub dark mode 下透明 SVG 看不清
  \begin{scope}[on background layer]
    \fill[white] (current bounding box.south west) rectangle (current bounding box.north east);
  \end{scope}
\end{tikzpicture}
\end{document}
